\section{转型期 CEO：品牌监护人、51/49 决策与创业精神}

前两章讲市场和技术，本章看组织和 CEO，因为转型不是技术路线图自己执行。这里要回答的问题是：当一家公司拥有 139 年历史、16 万员工、全球市场和巨大存量客户时，CEO 如何既保持进攻，又不摧毁品牌信用。康林松在 2019 年接任 CEO 时，已经在奔驰工作多年，但真正成为 CEO 后，“回头再看，没有其他人，只有你”。这个表达很关键：在 16 万员工、150 多个国家的大系统中，最终责任会落到 CEO 身上，而转型期的决策往往不是 80/20，而是 51/49。

\begin{figure}[H]
\centering
\includegraphics[width=0.96\textwidth]{transformation-ceo-loop.png}
\caption{转型期 CEO 决策闭环：51/49 决策、持续学习、必要时修正。自制概念图，依据 00:40:17--00:50:00 对谈内容整理。}
\end{figure}

\begin{knowledgebox}{读图：转型期决策不是一次押注}
图中 Uncertainty、Decision、Observe、Correct 和 Guardian 构成闭环。康林松强调，如果后来发现选择错了，就要收集新信息并修正，不要困于面子或沉没成本。
\end{knowledgebox}

\subsection{职业经理人如何保持创业精神}

上一节说明 CEO 的责任不可转移，本节看成熟公司如何避免组织钝化。张小珺问到，特斯拉和中国车企仍然创始人驱动，而德国汽车制造商经历多代职业经理人，是否更保守？康林松的回答是：成熟公司也必须保持创业精神，同时面对金融市场压力和转型压力。不能只打安全牌，也不能只防守。要投资新技术、提高效率、精简架构，让组织保持竞技状态。

\begin{warningbox}{大公司转型的双重压力}
初创公司可以用亏损换速度，成熟上市公司同时面对利润、股东、品牌、员工和转型投入。转型期 CEO 的难点在于既要保住财务实力，又不能因为安全牌错过技术窗口。
\end{warningbox}

\subsection{精简、速度和多车型战略的张力}

上一节讲成熟公司要保持创业精神，本节进入具体管理矛盾。康林松一边强调精简、速度和效率，一边又说未来三年奔驰会推出比以往更多的车型。这看似矛盾，实际是攻守关系：组织流程要精简，产品进攻不能停；成本要控制，但前锋也要上场。用他的足球比喻，只靠守门员无法赢球，成熟公司不能只靠防守保利润。

\noindent\begin{tabularx}{\textwidth}{p{0.20\textwidth}p{0.34\textwidth}X}
\toprule
管理目标 & 容易误解 & 更准确理解 \\
\midrule
精简 & 少做产品、减少投入 & 减少层级和流程浪费，提高决策速度。 \\
效率 & 只压成本 & 把研发、工厂、供应链和合作伙伴组织得更敏捷。 \\
多车型战略 & 与精简矛盾 & 产品进攻必须保留，组织复杂度要被管理。 \\
创业精神 & 像初创公司一样不顾利润 & 在财务约束内保持进攻和新技术投入。 \\
\bottomrule
\end{tabularx}

\subsection{品牌监护人}

前面讲组织如何行动，本节回到 CEO 如何定义成功。康林松对成功转型的定义，不是单一市场份额、利润率或技术领先，而是作为 Mercedes-Benz 品牌的监护人，在交给下一任时让品牌、客户体验、财务实力和创新实力都更好。这个定义很适合解释百年工业公司的时间尺度：技术会变，发动机会变，能源会变，但品牌承诺必须持续。

\begin{figure}[H]
\centering
\includegraphics[width=0.94\textwidth]{guardian-of-brand.png}
\caption{品牌监护人：成功不是单一技术指标，而是交接时品牌更强。自制概念图，依据 00:51:03--00:57:57 对谈内容整理。}
\end{figure}

\begin{knowledgebox}{读图：短期指标要服从长期监护}
左侧短期指标包括市场份额、利润率、技术领先和销量；右侧监护人目标包括品牌、客户体验、财务实力和创新实力。奔驰的 CEO 不能只优化一个季度，也不能只追一个技术标签。
\end{knowledgebox}

\subsection{不被任何一种技术绑定}

前面讲品牌监护人的指标，本节进入更底层的品牌身份。访谈最后，张小珺问汽油车是否会完全消失，以及如果那一天到来奔驰会是谁。康林松的回答是：奔驰正在走向零排放，但奔驰不会被任何一种特定技术绑定。技术总是在变，奔驰一直存在。这个回答很像百年品牌的底层原则：不要把身份绑死在发动机、电池、屏幕或某个软件功能上，而要把身份绑定在更稳定的承诺上。

\begin{importantbox}{技术会变，品牌承诺不能变}
对奔驰来说，燃油发动机、电动平台、AI 座舱和自动驾驶都是阶段性技术。品牌真正要守住的是安全、质量、工程、体验、美学和客户信任。转型的目标不是守住旧技术，而是让这些承诺穿过新技术周期。
\end{importantbox}

\subsection{转型风险清单}

本节把 CEO 视角压缩成风险清单。百年公司转型失败，往往不是因为看不到趋势，而是因为多个约束互相拉扯：怕损伤存量客户而转得太慢，怕错过新技术而投得太散，怕组织复杂而决策太慢，怕财务压力而减少创新。康林松的回答可以理解为在这些风险之间做动态平衡。

\noindent\begin{tabularx}{\textwidth}{p{0.20\textwidth}p{0.34\textwidth}X}
\toprule
风险 & 表现 & 对应治理方式 \\
\midrule
转型过慢 & 只守燃油和传统豪华 & 投入 MBOS、AI、电动平台和中国研发。 \\
转型过猛 & 伤害保值、品牌和利润 & 保留产品选择，维护高端客户价值。 \\
技术失控 & 外部伙伴技术碎片化 & 用 MBOS 做架构整合和体验控制。 \\
组织迟缓 & 大公司层级太多 & 精简流程，提高效率和创业精神。 \\
短期指标绑架 & 只追市场份额或季度利润 & 以品牌监护人身份做长期平衡。 \\
\bottomrule
\end{tabularx}

\subsection{本章小结}

转型期 CEO 的工作是把不确定性变成持续修正的组织能力。奔驰要回到创立时的开拓精神，同时保持百年品牌的长期承诺。

